\section{Task and Data}
\label{sec:task}

\paragraph{Taxonomy.} We use the taxonomy of \citet{cattan2025dragged}. Given a query and its retrieved sources,
the conflict is one of: \emph{No conflict} (the sources agree, up to wording or granularity);
\emph{Complementary information} (the answers differ but can all be correct, e.g.\ because the query is
underspecified); \emph{Conflicting opinions or research outcomes} (a subjective or contested topic with genuinely
opposing views); \emph{Outdated information} (a factual question on which sources disagree because the answer
changed over time); and \emph{Misinformation} (a factual question on which some sources are wrong).

\paragraph{Data.} \conflicts contains \nall queries drawn from ConflictingQA
\citep{wan2024evidence}, the geographic and temporal subsets of SituatedQA \citep{zhang2021situatedqa}, FreshQA
\citep{vu2023freshllms} and QACC \citep{liu2024qacc}, each with the top Google Search results
(\avgresults on average), the conflict type assigned by expert annotators, and, for the types with a single correct
answer, that answer. Each result carries a URL, title, snippet, publication date when available, and the
512-token segment of the page most relevant to the query. The released file has no split. We draw a stratified
development split of \ndev queries with a fixed seed and use it only to choose settings; the remaining \ntest
queries form the test split on which we report (Table~\ref{tab:data}). \emph{Misinformation} has only five
instances in the whole dataset, a limitation we return to in \S\ref{sec:diagnostics} by constructing a planted
variant.

\begin{table}[t]
\centering\small
\fitwidth{\begin{tabular}{lrrr}
\toprule
Conflict type & Dev & Test & All \\
\midrule
No conflict & 18 & 143 & 161 \\
Complementary information & 13 & 102 & 115 \\
Conflicting opinions & 13 & 102 & 115 \\
Outdated information & 7 & 55 & 62 \\
Misinformation & 1 & 4 & 5 \\
\midrule
Total & 52 & 406 & 458 \\
\bottomrule
\end{tabular}}
\caption{CONFLICTS \citep{cattan2025dragged} and our frozen split. The development split (stratified, fixed seed) is used only to choose settings; all reported results are on the test split unless stated.}
\label{tab:data}
\end{table}


\paragraph{Metric.} Following \citet{cattan2025dragged}, the primary metric is five-class accuracy. We also report
macro-F1 (the classes are imbalanced), per-class F1, and, for systems that return probabilities, expected
calibration error \citep[ECE;][]{guo2017calibration} of the top-class probability.
